\documentclass[11pt]{article}
\usepackage[letterpaper,margin=1in]{geometry}
\usepackage[T1]{fontenc}
\usepackage{lmodern}
\usepackage{amsmath,amssymb,amsthm}
\usepackage{microtype}
\usepackage[hidelinks]{hyperref}
\hypersetup{pdftitle={Quadratic convergence of SOS-density bounds on convex bodies},pdfauthor={Alec Kriebel}}
\newtheorem{theorem}{Theorem}
\newtheorem{lemma}{Lemma}
\newtheorem{proposition}{Proposition}
\newcommand{\R}{\mathbb R}
\newcommand{\vol}{\operatorname{vol}_n}
\newcommand{\op}{\mathrm{op}}
\newcommand{\dd}{\,\mathrm{d}}
\title{Quadratic convergence of SOS-density bounds\\on convex bodies}
\author{Alec Kriebel\\
  \small ORCID: \href{https://orcid.org/0009-0001-9320-500X}{0009-0001-9320-500X}}
\date{17 October 2026}

\begin{document}
\maketitle

\begin{abstract}
For every fixed compact convex body $K\subset\R^n$ with nonempty
interior and every real polynomial $f$, the normalized Lebesgue
sum-of-squares density bound of degree at most $2r$ has error
$O_{K,f}(r^{-2})$ for every integer $r\ge1$.
We give a self-contained proof and an explicit constant, using one
global square. A quadratic Taylor upper estimator reduces the objective
to a scalar endpoint problem. Homothetic contraction compares every
sublevel shell with the same core mass that normalizes a fixed power of
a Dirichlet polynomial; that mass cancels before summation. This proves
the quadratic upper rate in the convex-body/Lebesgue regime of the
2018 Oberwolfach question. A classical interval example shows that the
exponent is sharp in the worst case. The proxy and the broad
localization and relative-mass mechanisms have established precedents;
no firstness or new-method claim is made.
\end{abstract}

\section{Statement and source context}

All integrals over $K$ use $n$-dimensional Lebesgue measure.
Let $\Sigma_r$ be the set of polynomials that are sums of squares of
real polynomials of degree at most $r$. These are exactly the global
sum-of-squares polynomials of degree at most $2r$: the highest-degree
homogeneous squares cannot cancel. For a compact set $K$ of positive
volume and a real polynomial $f$, define
\begin{equation}\label{eq:hierarchy}
 H_r(f,K)=\inf\left\{\int_K f(x)q(x)\dd x:
       q\in\Sigma_r,\quad \int_Kq(x)\dd x=1\right\}.
\end{equation}
The constant density makes this set nonempty, and positivity gives
$H_r(f,K)\ge f_*:=\min_K f$.

\begin{theorem}\label{thm:main}
Let $n\ge1$, let $K\subset\R^n$ be compact and convex with nonempty
interior, and let $f\in\R[x_1,\ldots,x_n]$. Choose any $a\in K$ with
$f(a)=f_*$. Set
\begin{align*}
 H&=\max_{x\in K}\|\nabla^2f(x)\|_{\op}, &
 \gamma&=1+H/2,\\
 g(x)&=\nabla f(a)\cdot(x-a)+\gamma\|x-a\|^2, &
 G&=\max_{x\in K}g(x),\\
 L&=n+2, &
 C_n&=4L^2\left(\frac\pi2\right)^{2L}(1+2^{n+2}).
\end{align*}
Then $0<G<\infty$ and, for every integer $r\ge1$,
\begin{equation}\label{eq:rate}
 0\le H_r(f,K)-f_*\le \frac{C_nG}{r^2}.
\end{equation}
A single global square, normalized on the whole of $K$, witnesses
the upper bound.
\end{theorem}

No smoothness or semialgebraicity of the boundary is assumed, and $f$
need not be convex. The constant is deliberately coarse. If
$D=\operatorname{diam}(K)$ and $B=\max_K\|\nabla f\|$, then
$G\le BD+(1+H/2)D^2$, so a bound independent of the choice of minimizer
is available. The construction is an existence proof; it does not give
an efficient method for locating $a$.

Laurent's contribution, joint work with de Klerk, to the 2018
Oberwolfach report asks for the convergence regime for general convex
bodies with Lebesgue measure \cite{Laurent2018}. Its displayed density
constraints (3)--(4) omit the value of the integral. We explicitly use
unit mass, as required by its probability formulation and by the full
hierarchy definition \eqref{eq:hierarchy} in
Slot--Laurent \cite[Eq.~(3)]{SlotLaurent2020}. Thus our statement concerns
the normalized problem, rather than a literal unnormalized infimum.
Slot--Laurent's Theorem~11 gives
$O(\log^2 r/r^2)$ for this convex-body/Lebesgue hierarchy, and their
concluding questions ask whether the logarithm can be removed.
Theorem~\ref{thm:main} supplies the quadratic upper rate in that
precise setting.

The constrained-gradient and quadratic-upper-estimator reductions are
already present in de Klerk--Laurent
\cite[preprint Cor.~4.3]{deKlerkLaurent2018}
and Slot--Laurent \cite[Lemmas~1 and~4]{SlotLaurent2020}.
Endpoint polynomial concentration is classical. In particular,
Totik--Varga \cite[Thm.~2.1]{TotikVarga2012}, with endpoint parameter
$a=0$ and $\psi(v)=\sqrt v$, supplies degree-$m$ polynomials equal to
$1$ at zero, bounded by $2$ on $[0,1]$, and bounded by
$A\exp[-c(m^2t)^{1/4}]$ for $t>0$, with $A,c>0$ independent of $m$.
Markov's inequality gives a positive core of width proportional to
$m^{-2}$; the relative-mass estimate below then gives another deduction
of our theorem from those older ingredients. Totik--Varga also compare
shell masses with local core mass in their Christoffel-function
arguments. We assign no originality to the proxy, endpoint
localization, or broad shell/core cancellation mechanism.

Our proof instead uses an explicit finite trigonometric identity and
imports no approximation-theory convergence theorem. The accompanying
source audit distinguishes a deduction assembled here from an earlier
explicitly printed resolution. Its bounded search found no such
full-target theorem in the inspected passages; this does not establish
firstness, historical disclosure, or current worldwide openness.
The weaker rate recorded in the final 2026 overview
\cite[Table~1]{LaurentSlot2026} is evidence about that source's contents,
not an openness certificate.

\section{An elementary concentration proof}

\subsection*{Quadratic proxy and relative sublevel mass}
For $x\in K$, convexity places $a+s(x-a)$ in $K$ for $0\le s\le1$.
The right derivative at zero of $f(a+s(x-a))$ is nonnegative, so
\[
 \ell(x):=\nabla f(a)\cdot(x-a)\ge0.
\]
Taylor's formula on this segment gives
\[
 0\le f(x)-f(a)
 \le \ell(x)+\frac H2\|x-a\|^2\le g(x).
\]
The added positive quadratic ensures $g(x)>0$ for $x\in K\setminus\{a\}$.
Full dimension and compactness therefore give $0<G<\infty$.
The polynomial $u=g/G$ has degree at most two and takes values in
$[0,1]$ on $K$. For $0\le\alpha\le1$,
\begin{equation}\label{eq:contraction}
 g(a+\alpha(x-a))
   =\alpha\ell(x)+\alpha^2\gamma\|x-a\|^2
   \le\alpha g(x).
\end{equation}
For every $t>0$ put
$S_t=\{x\in K:u(x)\le t\}$ and $F(t)=\vol(S_t)$, so that
$F(t)=\vol(K)$ for $t\ge1$.
If $0<s\le t$, the homothety
$x\mapsto a+(s/t)(x-a)$ maps $S_t$ into $S_s$ by
\eqref{eq:contraction}. Its Jacobian is $(s/t)^n$, hence
\begin{equation}\label{eq:growth}
 F(t)\le(t/s)^n F(s)\qquad(0<s\le t).
\end{equation}
In particular $F(h)\ge h^n\vol(K)>0$ for $0<h\le1$.
The estimate also holds when $t>1$, as needed for an outer shell.
No density formula or differentiation of $F$ is used.

\subsection*{A polynomial endpoint kernel}
For an integer $m\ge1$, let $N=2m+1$ and define
\begin{equation}\label{eq:kernel}
 A_m(t)=\frac{1+2\sum_{j=1}^m T_j(1-2t)}{N},
\end{equation}
where $T_j(\cos\theta)=\cos(j\theta)$ is the Chebyshev polynomial.
Thus $A_m$ is a real polynomial of degree $m$ and $A_m(0)=1$.
Writing $v=\arcsin\sqrt t$ for $0<t\le1$, the finite Dirichlet-sum
identity gives
\begin{equation}\label{eq:sine}
 A_m(t)=\frac{\sin(Nv)}{N\sin v}
        =\frac{\sin(N\arcsin\sqrt t)}{N\sqrt t}.
\end{equation}
The defining sum and the sine quotient respectively imply
\begin{equation}\label{eq:envelope}
 |A_m(t)|\le\min\{1,(N\sqrt t)^{-1}\},\qquad 0<t\le1.
\end{equation}
Put $h=N^{-2}$. Convexity of arcsine gives
$\arcsin z\le\pi z/2$ on $[0,1]$.
For $0<t\le h$ we consequently have $Nv\le\pi/2$.
Using $\sin w\ge2w/\pi$ there, and $v\ge\sqrt t$, in
\eqref{eq:sine} yields
\begin{equation}\label{eq:core}
 A_m(t)\ge 2/\pi\qquad(0\le t\le h),
\end{equation}
including $t=0$ since $A_m(0)=1$.

\subsection*{Cancellation of the core mass}
Take $L=n+2$ and $Q_m(x)=A_m(u(x))^{2L}$.
This is the global square $[A_m(u(x))^L]^2$, regardless of its values
off $K$. Equations \eqref{eq:envelope}--\eqref{eq:core} give
\begin{equation}\label{eq:normalizer}
 Z_m:=\int_KQ_m\dd x\ge(2/\pi)^{2L}F(h)>0,
 \qquad Q_m\le\min\{1,(h/u)^L\}\quad(u>0).
\end{equation}
The core $\{u\le h\}$ contributes at most $hF(h)$ to
$\int_KuQ_m\dd x$. For the pairwise disjoint shells
\[
 E_j=\{x\in K:2^jh<u(x)\le2^{j+1}h\},\qquad j=0,1,\ldots,
\]
equation \eqref{eq:growth} gives
\begin{align*}
 \int_{E_j}uQ_m\dd x
 &\le 2^{j+1}h\,2^{-jL}F(2^{j+1}h)\\
 &\le hF(h)\,2^{n+1}2^{j(n+1-L)}.
\end{align*}
Only finitely many shells meet $K$. Extending their nonnegative
bounds to an infinite sum is legitimate, and $L=n+2$ gives
$\sum_{j\ge0}2^{j(n+1-L)}=2$. Thus
\[
 \int_KuQ_m\dd x\le hF(h)(1+2^{n+2}).
\]
Dividing by \eqref{eq:normalizer} cancels the same $F(h)$:
\begin{equation}\label{eq:ratio}
 \frac{\int_KuQ_m\dd x}{Z_m}
 \le \left(\frac\pi2\right)^{2L}
       \frac{1+2^{n+2}}{(2m+1)^2}.
\end{equation}
Zeros and oscillation of $A_m$ away from the core cause no difficulty;
the denominator uses its positive core and the tails use absolute
upper bounds over the whole body.

\subsection*{Degree and every integer order}
For $r\ge2L$ choose $m=\lfloor r/(2L)\rfloor\ge1$.
Since $\deg u\le2$,
\[
 \deg(A_m(u)^L)\le2mL\le r,\qquad \deg Q_m\le4mL\le2r.
\]
The density
\[
 q_r(x)=\frac{Q_m(x)}{Z_m}
       =\left[\frac{A_m(u(x))^L}{\sqrt{Z_m}}\right]^2
\]
is therefore feasible for \eqref{eq:hierarchy}. As $f-f(a)\le Gu$,
\eqref{eq:ratio} bounds its error. Also
$2\lfloor r/(2L)\rfloor+1\ge r/(2L)$, giving
\eqref{eq:rate} with the stated $C_n$.
For $1\le r<2L$, use $q_r=1/\vol(K)$, whose error is at most $G$.
Here $C_n/r^2\ge1$, so the same bound holds.
This proves Theorem~\ref{thm:main} for every integer $r\ge1$.
\hfill$\square$

\section{The known worst-case exponent}

Interval linear objectives already give matching quadratic lower
rates in de Klerk--Laurent
\cite[preprint Thm.~3.1 and Cors.~3.2--3.3]{deKlerkLaurent2020}.
The following elementary, coarse bound recalls why the universal
exponent cannot be improved.

\begin{proposition}\label{prop:lower}
For $K=[0,1]$ and $f(t)=t$,
\[
 H_r(f,K)\ge\frac1{4(r+1)^2}\qquad(r\ge1).
\]
The same lower bound holds for
$K=[0,1]\times[-1,1]^{n-1}$ and $f(t,y)=t$ in every fixed dimension
$n\ge1$.
\end{proposition}

\begin{proof}
Let $P_j$ be the Legendre polynomial normalized by $P_j(1)=1$.
The classical identities
\[
 \int_{-1}^1 P_j(x)P_k(x)\dd x=
       \frac{2\delta_{jk}}{2j+1},\qquad
 |P_j(x)|\le1\quad(-1\le x\le1)
\]
can be obtained directly from Rodrigues' formula
$P_j(x)=(2^jj!)^{-1}(\mathrm d/\mathrm dx)^j(x^2-1)^j$:
integration by parts gives orthogonality and the norm. For the
supremum bound, expand the finite-polynomial identity
\[
 P_j(x)=\frac1\pi\int_0^\pi
       \bigl(x+\iota\sqrt{1-x^2}\cos\theta\bigr)^j\dd\theta,
       \qquad \iota^2=-1;
\]
the base has modulus at most $1$.
Consequently $\ell_j(t)=\sqrt{2j+1}P_j(2t-1)$ is an orthonormal
basis on $[0,1]$ for polynomials of degree at most $r$ when
$0\le j\le r$. Cauchy--Schwarz gives, for each such polynomial $p$,
\[
 |p(t)|^2\le\left(\sum_{j=0}^r\ell_j(t)^2\right)
                 \int_0^1p(s)^2\dd s
             \le(r+1)^2\int_0^1p(s)^2\dd s.
\]
If $q=\sum_jp_j^2$ is feasible, summing shows $0\le q(t)\le M$,
where $M=(r+1)^2$, since its integral is $1$.
With $\alpha=1/(2M)$, we have $\int_0^\alpha q\le1/2$; hence
\[
 \int_0^1tq(t)\dd t\ge\alpha\int_\alpha^1q(t)\dd t
                \ge\frac1{4M}.
\]
Taking the infimum proves the interval assertion. On the product
body, apply the polynomial bound with $y$ fixed to every $p_j(t,y)$,
then sum and integrate over $y$. The resulting marginal density in
$t$ has integral $1$ and is bounded by $M$, so the same argument
applies.
\end{proof}

Together with Theorem~\ref{thm:main}, this gives one fixed instance
with error $\Theta(r^{-2})$ in each dimension. It does not give a
positive matching lower rate for each objective or each body; constant
objectives have zero error.

\section{Scope, supporting material, and review status}

The reference measure in Theorem~\ref{thm:main} is Lebesgue measure.
The homothety identity is not asserted for arbitrary reference
measures. The construction also uses a quadratic proxy $g$, rather
than densities constrained to be polynomials of the original
objective $f$; it therefore does not settle the same rate for that
more restrictive hierarchy for general nonconvex $f$.
The proof extends to $f$ of class $C^2$ on an open neighborhood of
$K$, since the same Taylor bound applies and the density still uses
only the quadratic proxy.

The accompanying verification package contains proof and source
audits, independent derivations and adversarial reviews, exact
arithmetic consistency checks, and an exact moment-matrix formulation
at the nonsimple apex of a square pyramid. These computations
and finite checks supplement the analytic proof; they do not certify
its universal quantifiers. Source retrieval, passages actually read,
and bounded literature-search limitations are recorded separately.

AI tools were used extensively in developing the solution,
derivations, source checks, verification, drafting, and adversarial
reviews. Independent reviewing agents were AI systems, not human
peer reviewers. This research note is unrefereed; it has not received
external human peer review and is not a proof-assistant or other
machine-formalized proof certificate. The author is responsible for
its mathematical claims and presentation.

\begin{thebibliography}{9}

\bibitem{Laurent2018}
M.~Laurent (joint work with E.~de Klerk),
Convergence analysis of Lasserre measure-based approximations for
polynomial optimization,
in \emph{Combinatorial Optimization}, Oberwolfach Reports
\textbf{15}(4), report 50/2018 (2018), 2991--2993.
\href{https://doi.org/10.4171/owr/2018/50}{doi:10.4171/owr/2018/50}.

\bibitem{SlotLaurent2020}
L.~Slot and M.~Laurent,
Improved convergence analysis of Lasserre's measure-based upper
bounds for polynomial minimization on compact sets,
\emph{Mathematical Programming} \textbf{193} (2022), 831--871.
\href{https://arxiv.org/abs/1905.08142v2}{arXiv:1905.08142v2} (2020).
\href{https://doi.org/10.1007/s10107-020-01468-3}{doi:10.1007/s10107-020-01468-3}.
Theorem and lemma numbers cited here refer to this preprint version.

\bibitem{deKlerkLaurent2018}
E.~de Klerk and M.~Laurent,
Comparison of Lasserre's measure-based bounds for polynomial
optimization to bounds obtained by simulated annealing,
\emph{Mathematics of Operations Research} \textbf{43}(4) (2018),
1317--1325.
\href{https://doi.org/10.1287/moor.2017.0906}{doi:10.1287/moor.2017.0906}.
\href{https://arxiv.org/abs/1703.00744v1}{arXiv:1703.00744v1}.

\bibitem{TotikVarga2012}
V.~Totik and T.~Varga,
Non-symmetric fast decreasing polynomials and applications,
\emph{Journal of Mathematical Analysis and Applications}
\textbf{394} (2012), 378--390.
\href{https://doi.org/10.1016/j.jmaa.2012.03.006}{doi:10.1016/j.jmaa.2012.03.006}.

\bibitem{LaurentSlot2026}
M.~Laurent and L.~Slot,
An overview of convergence rates for sum of squares hierarchies in
polynomial optimization,
in \emph{Recent Developments in Industrial and Applied Mathematics},
ICIAM2023 Springer Series \textbf{1} (2026), 149--176.
\href{https://doi.org/10.1007/978-981-95-1446-5_8}{doi:10.1007/978-981-95-1446-5\_8}.
\href{https://ir.cwi.nl/pub/36052/36052.pdf}{Author-hosted final text}.

\bibitem{deKlerkLaurent2020}
E.~de Klerk and M.~Laurent,
Worst-case examples for Lasserre's measure-based hierarchy for
polynomial optimization on the hypercube,
\emph{Mathematics of Operations Research} \textbf{45}(1) (2020), 86--98.
\href{https://doi.org/10.1287/moor.2018.0983}{doi:10.1287/moor.2018.0983}.
\href{https://arxiv.org/abs/1804.05524v1}{arXiv:1804.05524v1}.

\end{thebibliography}
\end{document}
